\textbf{Phần 1. Trắc nghiệm lựa chọn}
\setcounterpageref{ex}{0}

\begin{ex} %[2D11V57]
Để đo chiều dài của một cuốn sách giáo khoa Toán 2, em nên dùng thước có vạch chia theo đơn vị nào là thích hợp nhất?
\choice
{m}
{km}
{dm}
{\True cm}
\loigiai{
Đo chiều dài sách dùng xăng-ti-mét (cm) là thích hợp nhất.
\textbf{Đáp án của câu hỏi là D}
}
\end{ex}

\begin{ex} %[2D11V57]
Một gang tay của học sinh lớp 2 dài khoảng bao nhiêu?
\choice
{15 m}
{\True 15 cm}
{15 dm}
{15 km}
\loigiai{
Một gang tay học sinh dài khoảng 15 cm.
\textbf{Đáp án của câu hỏi là B}
}
\end{ex}

\begin{ex} %[2D11V57]
Chiều cao của cánh cửa lớp học thường khoảng
\choice
{2 cm}
{20 m}
{\True 20 dm}
{20 cm}
\loigiai{
Cánh cửa lớp học cao khoảng 2 m = 20 dm.
\textbf{Đáp án của câu hỏi là C}
}
\end{ex}

\begin{ex} %[2D11V57]
Một sợi dây khi gập đôi lại thì đo được $4\text{ dm}$. Chiều dài sợi dây đó khi duỗi thẳng là
\choice
{\True 8 dm}
{2 dm}
{6 dm}
{10 dm}
\loigiai{
$4 \times 2 = 8$ dm.
\textbf{Đáp án của câu hỏi là A}
}
\end{ex}

\begin{ex} %[2D11V57]
Dùng thước kẻ dài $20\text{ cm}$ để đo mép bàn. Đặt thước liên tiếp $5$ lần thì vừa hết chiều dài mép bàn. Chiều dài mép bàn đó là
\choice
{10 cm}
{50 cm}
{2 m}
{\True 1 m}
\loigiai{
$20 \times 5 = 100\text{ cm} = 1\text{ m}$.
\textbf{Đáp án của câu hỏi là D}
}
\end{ex}

\begin{ex} %[2D11V57]
Khoảng cách giữa hai cột điện liên tiếp trên phố thường được đo bằng đơn vị nào?
\choice
{cm}
{\True m}
{mm}
{dm}
\loigiai{
Khoảng cách giữa 2 cột điện đo bằng mét (m).
\textbf{Đáp án của câu hỏi là B}
}
\end{ex}

\begin{ex} %[2D11V57]
Để đo chu vi của thân một cây đại thụ trong sân trường, cách làm nào sau đây thuận tiện nhất?
\choice
{Dùng thước kẻ thẳng nhựa ngắn 15 cm}
{Dùng cân đĩa}
{\True Dùng một sợi dây mềm quấn một vòng quanh thân cây rồi đo chiều dài đoạn dây đó}
{Dùng đồng hồ bấm giây}
\loigiai{
Dùng sợi dây mềm quấn quanh thân cây rồi đo đoạn dây là thuận tiện và chính xác nhất.
\textbf{Đáp án của câu hỏi là C}
}
\end{ex}

\begin{ex} %[2D11V57]
Một mảnh vườn hình vuông có cạnh dài $25\text{ m}$. Bác thợ rào xung quanh mảnh vườn bằng lưới thép. Chiều dài hàng rào là
\choice
{\True 100 m}
{50 m}
{75 m}
{125 m}
\loigiai{
Chu vi hình vuông: $25 \times 4 = 100$ m.
\textbf{Đáp án của câu hỏi là A}
}
\end{ex}

\begin{ex} %[2D11V57]
Một bạn bước chân được khoảng $4\text{ dm}$ mỗi bước. Bạn đi từ cửa lớp đến bảng hết $10$ bước chân. Khoảng cách đó dài khoảng bao nhiêu mét?
\choice
{40 m}
{2 m}
{14 m}
{\True 4 m}
\loigiai{
$4 \times 10 = 40\text{ dm} = 4\text{ m}$.
\textbf{Đáp án của câu hỏi là D}
}
\end{ex}

\begin{ex} %[2D11V57]
Một cuộn băng dính dài $5\text{ m}$. Người ta cắt $5$ đoạn để dán trang trí, mỗi đoạn dài $4\text{ dm}$. Cuộn băng dính còn lại bao nhiêu đề-xi-mét?
\choice
{20 dm}
{\True 30 dm}
{10 dm}
{25 dm}
\loigiai{
Đổi $5\text{ m} = 50\text{ dm}$. Đã cắt: $4 \times 5 = 20$ dm. Còn lại: $50 - 20 = 30$ dm.
\textbf{Đáp án của câu hỏi là B}
}
\end{ex}

\begin{ex} %[2D11V57]
Tổ Một đo được đoạn đường dài $12\text{ m}$, tổ Hai đo được đoạn đường dài $15\text{ m}$. Cả hai tổ đo được tổng cộng bao nhiêu mét đường?
\choice
{25 m}
{26 m}
{\True 27 m}
{28 m}
\loigiai{
$12 + 15 = 27$ m.
\textbf{Đáp án của câu hỏi là C}
}
\end{ex}

\begin{ex} %[2D11V57]
Một thanh tre dài $2\text{ m } 4\text{ dm}$. Bác thợ muốn cắt thành $4$ đoạn bằng nhau. Mỗi đoạn dài bao nhiêu đề-xi-mét?
\choice
{\True 6 dm}
{5 dm}
{4 dm}
{7 dm}
\loigiai{
$2\text{ m } 4\text{ dm} = 24\text{ dm}$. Mỗi đoạn: $24 : 4 = 6$ dm.
\textbf{Đáp án của câu hỏi là A}
}
\end{ex}

\textbf{Phần 2. Trắc nghiệm đúng sai}
\setcounterpageref{ex}{0}

\begin{ex} %[2D11V57]
Xét tính hợp lý của các ước lượng và kết quả đo độ dài thực tế:
\choiceTF
{\True Một bước chân người lớn dài khoảng 70 cm.}
{Độ dày của một quyển vở ô ly khoảng 5 dm.}
{\True Khi gập đôi một sợi dây 3 lần, chiều dài mỗi phần bằng một phần tám sợi dây ban đầu.}
{\True 100 bước chân mỗi bước 5 dm dài tất cả 50 m.}
\loigiai{
\begin{itemize}
    \item Ý a) Đúng: Ước lượng phù hợp thực tế. Đúng.
    \item Ý b) Sai: 5 dm = 50 cm là quá dày, quyển vở chỉ dày vài milimét.
    \item Ý c) Đúng: Gập 3 lần: chia làm 2 x 2 x 2 = 8 phần bằng nhau. Đúng.
    \item Ý d) Đúng: 100 x 5 dm = 500 dm = 50 m. Đúng.
\end{itemize}
\begin{itemize}
    \item a) \textbf{Đúng}.
    \item b) \textbf{Sai}.
    \item c) \textbf{Đúng}.
    \item d) \textbf{Đúng}.
\end{itemize}
}
\end{ex}

\begin{ex} %[2D11V57]
Trong tiết trải nghiệm, Lan dùng thước cuộn đo kích thước phòng học:
\choiceTF
{\True Chiều dài phòng học đo được 8 m.}
{\True Chiều rộng phòng học đo được 6 m.}
{Chu vi sàn phòng học là 14 m.}
{\True Nửa chu vi sàn phòng học là 14 m.}
\loigiai{
\begin{itemize}
    \item Ý a) Đúng: Kích thước thực tế hợp lý. Đúng.
    \item Ý b) Đúng: Kích thước thực tế hợp lý. Đúng.
    \item Ý c) Sai: Chu vi: (8 + 6) x 2 = 28 m, không phải 14 m.
    \item Ý d) Đúng: 8 + 6 = 14 m. Đúng.
\end{itemize}
\begin{itemize}
    \item a) \textbf{Đúng}.
    \item b) \textbf{Đúng}.
    \item c) \textbf{Sai}.
    \item d) \textbf{Đúng}.
\end{itemize}
}
\end{ex}

\textbf{Phần 3. Trả lời ngắn}
\setcounterpageref{ex}{0}

\begin{ex} %[2D11V57]
Dùng thước thẳng 30 cm đo mép bàn. Đặt thước 4 lần liên tiếp thì còn thiếu 10 cm nữa mới hết mép bàn. Chiếc bàn đó dài bao nhiêu đề-xi-mét?
\shortans{13}
\loigiai{
\textbf{Đáp số:} $13$\\
$30 \times 4 + 10 = 130\text{ cm} = 13$ dm.
}
\end{ex}

\begin{ex} %[2D11V57]
Một đoạn dây gập làm 4 phần bằng nhau. Một phần dài 25 cm. Cả sợi dây dài bao nhiêu mét?
\shortans{1}
\loigiai{
\textbf{Đáp số:} $1$\\
$25 \times 4 = 100\text{ cm} = 1$ m.
}
\end{ex}

\begin{ex} %[2D11V57]
Bạn Hùng bước 50 bước, mỗi bước khoảng 4 dm. Sân trường dài khoảng bao nhiêu mét?
\shortans{20}
\loigiai{
\textbf{Đáp số:} $20$\\
$50 \times 4 = 200\text{ dm} = 20$ m.
}
\end{ex}

\begin{ex} %[2D11V57]
Bác An có sợi dây thừng dài 15 m. Bác cắt làm 3 đoạn bằng nhau. Mỗi đoạn dài bao nhiêu mét?
\shortans{5}
\loigiai{
\textbf{Đáp số:} $5$\\
$15 : 3 = 5$ m.
}
\end{ex}

\begin{ex} %[2D11V57]
Một hình tam giác có độ dài ba cạnh lần lượt là 12 cm, 18 cm, 20 cm. Chu vi của hình tam giác đó là bao nhiêu đề-xi-mét?
\shortans{5}
\loigiai{
\textbf{Đáp số:} $5$\\
$12 + 18 + 20 = 50\text{ cm} = 5$ dm.
}
\end{ex}

\begin{ex} %[2D11V57]
Một con kiến bò quanh mép một hộp bút hình chữ nhật dài 15 cm, rộng 5 cm. Kiến bò hết một vòng là bao nhiêu xăng-ti-mét?
\shortans{40}
\loigiai{
\textbf{Đáp số:} $40$\\
$(15 + 5) \times 2 = 40$ cm.
}
\end{ex}
